\section{Билет 40. Основное уравнение молекулярно-кинетической теории. Постоянная Больцмана. Средняя кинетическая энергия поступательного движения одноатомной молекулы и ее связь с температурой. Среднеквадратичная скорость молекулы газа.}

\begin{definition}
\textbf{Основное уравнение МКТ} связывает давление идеального газа с концентрацией молекул и их средней кинетической энергией:
\begin{equation}
    p = \frac{2}{3} n \langle E_{\text{пост}} \rangle, \label{eq:mkt_main_40}
\end{equation}
где $n = N/V$ --- концентрация молекул, $\langle E_{\text{пост}} \rangle = \langle m_0 v^2 / 2 \rangle$ --- средняя кинетическая энергия поступательного движения одной молекулы.
\end{definition}

\begin{theorem}[Связь энергии и температуры]
Сравнивая основное уравнение МКТ \eqref{eq:mkt_main_40} с уравнением состояния идеального газа, записанным через концентрацию:
\begin{equation}
    p = n k T, \label{eq:ideal_concentration}
\end{equation}
где $k$ --- постоянная Больцмана, получаем фундаментальное соотношение:
\begin{equation}
    \frac{2}{3} n \langle E_{\text{пост}} \rangle = n k T \quad \Rightarrow \quad \langle E_{\text{пост}} \rangle = \frac{3}{2} k T. \label{eq:energy_temp}
\end{equation}
\end{theorem}

\begin{property}
Из формулы \eqref{eq:energy_temp} следуют важные выводы:
\begin{enumerate}
    \item \textbf{Физический смысл температуры:} Абсолютная температура $T$ является мерой средней кинетической энергии поступательного движения молекул идеального газа. Чем выше температура, тем интенсивнее тепловое движение частиц.
    \item Температура --- \textbf{статистическая величина}: она характеризует не отдельную молекулу, а ансамбль из огромного числа частиц. Понятие «температура одной молекулы» лишено физического смысла.
    \item Для \textbf{одноатомных молекул} (инертные газы: He, Ne, Ar) вся внутренняя энергия сосредоточена в поступательном движении, поэтому формула \eqref{eq:energy_temp} описывает полную среднюю энергию одной молекулы.
    \item Средняя энергия поступательного движения \textbf{не зависит} от массы молекулы, химической природы газа и давления --- она определяется только температурой.
\end{enumerate}
\end{property}

\begin{definition}
\textbf{Среднеквадратичная скорость} молекул газа $v_{\text{кв}}$ --- корень квадратный из среднего значения квадрата скорости:
\begin{equation}
    v_{\text{кв}} = \sqrt{\langle v^2 \rangle}.
\end{equation}
Используя связь $\langle E_{\text{пост}} \rangle = m_0 \langle v^2 \rangle / 2 = 3kT/2$, получаем формулу для расчёта:
\begin{equation}
    v_{\text{кв}} = \sqrt{\frac{3kT}{m_0}} = \sqrt{\frac{3RT}{M}}, \label{eq:v_rms}
\end{equation}
где $M = N_A m_0$ --- молярная масса газа.
\end{definition}

\begin{remark}
\textbf{Численная оценка для воздуха.} При комнатной температуре ($T = 300$ К) и молярной массе воздуха $M \approx 0.029$ кг/моль:
\begin{equation}
    v_{\text{кв}} = \sqrt{\frac{3 \cdot 8.31 \cdot 300}{0.029}} \approx 508 \text{ м/с}.
\end{equation}
Это значение сравнимо со скоростью звука в воздухе ($\approx 340$ м/с), что подтверждает молекулярную природу звуковых волн.
\par\noindent\textbf{Важно:} Среднеквадратичная скорость $v_{\text{кв}}$ отличается от средней скорости $\langle v \rangle$ и наиболее вероятной скорости $v_{\text{н}}$. Для максвелловского распределения: $v_{\text{н}} : \langle v \rangle : v_{\text{кв}} \approx 1 : 1.13 : 1.22$.
\end{remark}

\begin{remark}
\textbf{Обобщение на многоатомные газы.} Для двухатомных и многоатомных молекул полная средняя энергия включает также вращательные и колебательные степени свободы. Согласно теореме о равнораспределении энергии, на каждую степень свободы приходится средняя энергия $kT/2$. Для одноатомного газа (3 поступательные степени свободы) получаем $\langle E \rangle = 3 \cdot (kT/2) = 3kT/2$, что согласуется с \eqref{eq:energy_temp}.
\end{remark}
